\section*{الفصل \ref{ch:g12:polymers} --- البوليمرات}

\begin{solution}{exo:g12:polymers:1}
البروبين، \ce{CH2=CH-CH3}.
\end{solution}

\begin{solution}{exo:g12:polymers:2}
\ce{-CF2-CF2-}؛ \omterm{def:g12:polymers:addition-polymer}{بوليمر إضافة}: فللمونومر رابطة \ce{C=C}
ولا يتكوّن أي \omterm{def:g7:chemical-reactions:reactant}{ناتج} آخر.
\end{solution}

\begin{solution}{exo:g12:polymers:3}
إضافة: البولي إيثين، وPVC، والبوليستيرين. تكاثف: PET (بوليستر)،
والنايلون (بولي \omterm{def:g11:functional-groups:amine}{أميد}).
\end{solution}

\begin{solution}{exo:g12:polymers:4}
السليلوز، والنشا، ومطاط اللاتكس، والبروتينات. أما النايلون وPVC
فاصطناعيان.
\end{solution}

\begin{solution}{exo:g12:polymers:5}
كل رابطة تستعمل مجموعة واحدة من كل \omterm{def:g12:polymers:polymer}{مونومر}؛ ومع مجموعتين، يستطيع
\omterm{def:g12:polymers:polymer}{المونومر} أن يرتبط من جانبيه فتستمر السلسلة في النمو. أما \omterm{def:g12:polymers:polymer}{المونومر} ذو
المجموعة الواحدة فيوقف السلسلة.
\end{solution}

\begin{solution}{exo:g12:polymers:6}
$M(\ce{C2H3Cl}) = 24.0 + 3.0 + 35.5 = \qty{62.5}{g/mol}$;
$n = 125000/62.5 = 2000$.
\end{solution}

\begin{solution}{exo:g12:polymers:7}
$M(\ce{C8H8}) = \qty{104.0}{g/mol}$; $M = 2500 \times 104.0 =
\qty{2.60e5}{g/mol}$.
\end{solution}

\begin{solution}{exo:g12:polymers:8}
المجموعة الحمضية \omterm{def:g7:atoms-and-molecules:molecule}{لجزيء} تكوّن إسترًا مع المجموعة الكحولية
\omterm{def:g7:atoms-and-molecules:molecule}{للجزيء} التالي، الذي تبقى له مجموعة حمضية حرة: فتنمو السلسلة. \omterm{def:g12:polymers:repeat-unit}{والوحدة
المتكررة} \ce{-O-CH(CH3)-CO-}؛ وينطلق الماء عند كل رابطة.
\end{solution}

\begin{solution}{exo:g12:polymers:9}
أغشية الأكياس مصنوعة من سلاسل متفرعة، لا يمكن رصّها: فالمادة
لا بلورية، لينة وشفافة. وقوارير الحليب مصنوعة من سلاسل شبه
خطية، تُرصّ في مناطق بلورية: فهي أكثف وأصلب،
والمناطق البلورية تبعثر الضوء.
\end{solution}

\begin{solution}{exo:g12:polymers:10}
سلاسلهما مرتبطة بجسور تساهمية كثيرة. فالتسخين لا يستطيع أن
يجعل السلاسل تنزلق؛ بل يكسر الروابط ويتلف المادة بدل
أن يصهرها.
\end{solution}

\begin{solution}{exo:g12:polymers:11}
$M(\ce{C12H22N2O2}) = 144.0 + 22.0 + 28.0 + 32.0 = \qty{226.0}{g/mol}$;
$n = 22600/226.0 = 100$.
\end{solution}

\begin{solution}{exo:g12:polymers:12}
\[
  \ce{H2N(CH2)6NH2 + HOOC(CH2)4COOH -> H2N(CH2)6NHCO(CH2)4COOH + H2O} .
\]
يبقى \omterm{def:g7:chemical-reactions:reactant}{للناتج} مجموعة \omterm{def:g11:functional-groups:amine}{أمين} حرة في أحد طرفيه ومجموعة حمضية حرة
في الطرف الآخر: فيمكن لكل منهما أن تتفاعل مع \omterm{def:g12:polymers:polymer}{مونومر} آخر.
\end{solution}

\begin{solution}{exo:g12:polymers:13}
جزيئا ماء لكل \omterm{def:g12:polymers:repeat-unit}{وحدة متكررة}. $n = 1000/226.0 =
\qty{4.42}{mol}$ من الوحدات المتكررة، إذن \qty{8.85}{mol} من الماء:
$8.85 \times 18.0 = \qty{159}{g}$.
\end{solution}

\begin{solution}{exo:g12:polymers:14}
المطاط الطبيعي الدافئ مكوّن من سلاسل ينزلق بعضها على بعض.
والجسور الكبريتية تربط السلاسل: فتتمدد لكنها ترتد، ولا
تعود تسيل. ولأن الإطار مترابط بالجسور، فلا يمكن صهره وإعادة تشكيله.
\end{solution}

\begin{solution}{exo:g12:polymers:15}
$M(\ce{C6H10O5}) = \qty{162.0}{g/mol}$؛ $n = \num{1.6e6}/162.0 = 9900$
وحدة؛ والطول $9900 \times 0.5 = \qty{4900}{nm}$، أي نحو
\qty{5}{\micro m}.
\end{solution}

\begin{solution}{pb:g12:polymers:1}
\textbf{1.} الإيثان-1,2-ديول (مجموعتا \omterm{def:g11:functional-groups:alcohol}{كحول})
وحمض البنزين-1,4-ثنائي الكربوكسيليك (مجموعتان حمضيتان كربوكسيليتان).

\textbf{2.} رابطة \omterm{def:g11:functional-groups:carboxylic-acid}{إستر}؛ وينطلق الماء.

\textbf{3.} \omterm{def:g12:polymers:addition-polymer}{بوليمر تكاثف}.

\textbf{4.} $M(\ce{C2H6O2}) = \qty{62.0}{g/mol}$؛ $M(\ce{C8H6O4}) =
\qty{166.0}{g/mol}$.

\textbf{5.} $62.0 + 166.0 - 2 \times 18.0 = \qty{192.0}{g/mol}$؛
و$10 \times 12.0 + 8 \times 1.0 + 4 \times 16.0 = 192.0$.

\textbf{6.} $n = 38400/192.0 = 200$.

\textbf{7.} رابطتا \omterm{def:g11:functional-groups:carboxylic-acid}{إستر} لكل \omterm{def:g12:polymers:repeat-unit}{وحدة متكررة}: نحو 400.

\textbf{8.} سلاسله خطية في معظمها، وتصطف أجزاء منها في
مناطق منتظمة.

\textbf{9.} إنه لدن حراري: فسلاسله غير مترابطة بجسور
وتنزلق بعضها على بعض حين تكون ساخنة.

\textbf{10.} $300/0.80 = \qty{375}{g}$.

\textbf{11.} $375/25 = 15$ قارورة.

\textbf{12.} تُنزع السدادات واللصاقات والأوساخ والقصاصات أو تضيع خلال
الفرز والغسل والغزل.

\textbf{13.} منع النفايات (المبدأ 1): فالشيء المستعمل يصير \omterm{def:g5:raw-materials-and-recycling:raw-material}{مادة
أولية}.

\textbf{14.} \ce{C10H8O4 + 2H2O -> C2H6O2 + C8H6O4}.

\textbf{15.} $1000/192.0 = \qty{5.21}{mol}$.

\textbf{16.} الديول: $5.21 \times 62.0 = \qty{323}{g}$؛ والحمض:
$5.21 \times 166.0 = \qty{865}{g}$.

\textbf{17.} $2 \times 5.21 \times 18.0 = \qty{188}{g}$ من الماء؛
$1000 + 188 = 323 + 865 = \qty{1188}{g}$: فالكتلة محفوظة.

\textbf{18.} يمكن تنقية المونومرات فتعطي PET جديدًا بجودة
الجديد، بينما يقصّر الصهر السلاسل قليلًا في كل مرة.

\textbf{19.} تحتاج السترة الصوفية الصناعية الواحدة إلى 15 قارورة كتلة كل منها \qty{25}{g}.
\end{solution}
